\documentclass[main.tex]{subfiles}
\begin{document}
\ifSubfilesClassLoaded{\begin{center}\small\itshape Part of the GnssProcessor Mathematical Reference \textbullet\ \today\end{center}}{}
\section{Inertial aiding}\label{sec:imu}
% ============================================================
%  Inertial aiding (tightly-coupled IMU integration)
%
%  Strapdown mechanization in the ECEF frame + 15-state INS
%  error model and how the GNSS update corrects it.
%
%  Scope note: this section documents ONLY the implemented core
%  (filter/ImuIntegrator, KalmanFilter::predictAided, the TC
%  Doppler Jacobian).  Lever-arm/scale-factor states and a full
%  Somigliana gravity model are NOT implemented today and are
%  called out as such inline.
% ============================================================

With an IMU configured, a first-order strapdown mechanization in the
Earth-centred Earth-fixed (ECEF) frame replaces the inter-epoch time
update of \dcref{sec:kf-predict}{the time update}.  Two switches control
it.  \texttt{imu.enabled} enables the IMU.  \texttt{estimation.aiding.imu}
selects the IMU source in the \texttt{AidingSourceManager}.  When
\texttt{estimation.aiding.imu} is absent it defaults to the value of
\texttt{imu.enabled}; an explicit value takes precedence
(\texttt{src/core/config/GnssConfigLoad.cpp}).

The mechanization is implemented in
\texttt{src/core/filter/ImuIntegrator.cpp}.  The per-epoch prediction
runs \texttt{Application::run} $\to$
\texttt{AidingSourceManager::collectPredict} $\to$
\texttt{src/core/aiding/ImuAidingSource.cpp}, and
\texttt{KalmanFilter::predictAided}
(\texttt{src/core/filter/KalmanFilter.cpp}) splices the resulting
\texttt{PredictContribution} into the filter covariance.  The real-time
engine (\texttt{RealTimeEngine::processEpoch}) uses the same
\texttt{predictAided} $/$ \texttt{collectPredict} path.  Only the sample
ingest differs.  Post-processing reads a CSV file with
\texttt{loadData}, which sorts the samples by time.  In real time each
sensor sample reaches \texttt{ImuIntegrator::addSample}, which appends it
to the same buffer in arrival order.  For samples that arrive in time
order the two paths therefore give the same \texttt{integrate} results.

The \texttt{IAidingSource} / \texttt{AidingSourceManager} abstraction
(\texttt{include/aiding/}) is the interface through which every non-GNSS
modality (IMU, ZUPT, NHC, baro, wheel odometry, vision) contributes up to
three items: a \texttt{StateManager}-owned state block, a predict-time
$\Phi/Q$ sub-block, and either independent measurement rows or extra
Jacobian \emph{columns} on existing GNSS rows.  The IMU is its first
client.  \texttt{ImuAidingSource} reproduces the arithmetic of the older
monolithic \texttt{KalmanFilter::predictImu}, so the
strapdown equations below apply unchanged.  \texttt{predictImu} remains
in the source but has no callers; \texttt{predictAided} replaces it.

The formulation is an \emph{error-state} extended Kalman filter.  The
mechanization propagates the nominal navigation state
$(\rxpos,\imuvel,\Cbe)$ deterministically from the inertial samples,
while a $15$-element INS error vector
\begin{equation}
\delta\stateV_{\mathrm{ins}}
  = \begin{bmatrix}\delta\rxpos & \delta\imuvel & \imupsi
    & \delta\biasa & \delta\biasg\end{bmatrix}^{\Tr}
  \in \mathbb{R}^{15}
\label{eq:ins-error-state}
\end{equation}
carries the linearised transition and process noise that propagate the
covariance.  Three attitude-error states $\imupsi$, three
accelerometer-bias states $\biasa$ and three gyro-bias states $\biasg$
augment the position state (and the three velocity states when
\texttt{estimation.estimate\_velocity} is set).  \texttt{StateManager}
allocates the nine added states as one contiguous block (accessors
\texttt{getAttitudeIndex}, \texttt{getAccelBiasIndex},
\texttt{getGyroBiasIndex}).

\begin{remark}[Implemented scope]
\texttt{ImuIntegrator} fixes the block order of
\cref{eq:ins-error-state}
($[\delta\rxpos,\delta\imuvel,\imupsi,\delta\biasa,\delta\biasg]$).  The
filter carries no lever-arm state and no sensor scale-factor states, and
the GNSS measurement model has no IMU-to-antenna lever-arm term
(\Cref{sec:imu-coupling}).  These are planned extensions.  The equations
below omit the corresponding terms.
\end{remark}

\subsection{Strapdown mechanization}\label{sec:imu-mech}

For each IMU sample, the raw specific force $\tilde{\boldsymbol a}$ and
angular rate $\tilde{\boldsymbol\omega}$ (both in the body frame) are
corrected with the current filter bias estimates,
\begin{equation}
\imuf = \tilde{\boldsymbol a} - \biasa, \qquad
\imuw = \tilde{\boldsymbol\omega} - \biasg ,
\label{eq:imu-bias-correct}
\end{equation}
and integrated over the sample interval $\dt$.  A sample whose interval
is not positive or exceeds \SI{10}{\second} is skipped.  The attitude
direction-cosine matrix $\Cbe$ (body $\to$ ECEF) is updated to first
order with Earth-rotation compensation,
\begin{empheq}[box=\fbox]{equation}
\Cbe \;\leftarrow\;
  \Cbe\bigl(\boldsymbol I + \skewm{\imuw}\,\dt\bigr)
  \;-\; \skewm{\earthrate}\,\Cbe\,\dt ,
\label{eq:imu-attitude}
\end{empheq}
where $\earthrate = [\,0,\,0,\,\omega_e\,]^{\Tr}$ is the ECEF
Earth-rotation vector
($\omega_e = \SI{7.2921151467e-5}{\radian\per\second}$,
\texttt{common/Constants.hpp}) and $\skewm{\cdot}$ is the skew-symmetric
(cross-product) matrix.  To limit numerical drift, $\Cbe$ is
re-orthonormalised by singular-value decomposition every $100$ integrated
samples, starting with the first.  The velocity and position propagate
as
\begin{empheq}[box=\fbox]{align}
\imuvel_{k+1} &= \imuvel_k
  + \bigl(\Cbe\,\imuf + \gravec(\rxpos_k)
          - 2\,\skewm{\earthrate}\,\imuvel_k\bigr)\dt ,
\label{eq:imu-velocity}\\[2pt]
\rxpos_{k+1}  &= \rxpos_k
  + \tfrac12\bigl(\imuvel_k + \imuvel_{k+1}\bigr)\dt .
\label{eq:imu-position}
\end{empheq}
The specific force is resolved with the attitude already updated by
\cref{eq:imu-attitude}, and gravity is evaluated at the position of the
start of the step.  \Cref{eq:imu-velocity} applies the Coriolis
acceleration $-2\,\skewm{\earthrate}\imuvel$.  The centrifugal
acceleration is part of the gravity vector $\gravec$: $\gravec$ is a
\emph{plumb-bob} gravity (gravitational plus centrifugal).
\Cref{eq:imu-position} is a trapezoidal position update.

The gravity vector is a point-mass radial model plus the centrifugal
term (\texttt{ImuIntegrator::computeGravityEcef}):
\begin{equation}
\gravec(\rxpos) =
  -\,\frac{\mu}{\lVert\rxpos\rVert^{3}}\,\rxpos
  \;+\; \omega_e^{2}\,[\,p_x,\,p_y,\,0\,]^{\Tr},
\qquad \mu = \SI{3.986005e14}{\cubic\metre\per\square\second}.
\label{eq:imu-gravity}
\end{equation}
The normal-gravity (Somigliana) flattening and the $J_2$ zonal term are
not modelled.  Over one GNSS epoch interval the resulting deterministic
acceleration error is small, and the estimated accelerometer bias
$\biasa$ absorbs it.

\subsection{INS error-state dynamics}\label{sec:imu-error}

Linearising \crefrange{eq:imu-attitude}{eq:imu-position} about the
nominal trajectory yields the continuous error dynamics
$\dot{\delta\stateV}_{\mathrm{ins}} = \insF\,\delta\stateV_{\mathrm{ins}}
+ \insG\,\boldsymbol w$, with, in the block order of
\cref{eq:ins-error-state},
\begin{equation}
\insF =
\begin{bmatrix}
\boldsymbol 0 & \boldsymbol I & \boldsymbol 0 & \boldsymbol 0 & \boldsymbol 0\\
\boldsymbol 0 & -2\skewm{\earthrate} & -\skewm{\imufe} & \Cbe & \boldsymbol 0\\
\boldsymbol 0 & \boldsymbol 0 & -\skewm{\earthrate} & \boldsymbol 0 & -\Cbe\\
\boldsymbol 0 & \boldsymbol 0 & \boldsymbol 0 & \boldsymbol 0 & \boldsymbol 0\\
\boldsymbol 0 & \boldsymbol 0 & \boldsymbol 0 & \boldsymbol 0 & \boldsymbol 0
\end{bmatrix},
\qquad \imufe := \Cbe\,\imuf .
\label{eq:ins-F}
\end{equation}
Both sensor biases are random walks (no deterministic dynamics: the
null bottom rows of \cref{eq:ins-F}).  Two terms of the standard ECEF
error model are omitted by design: the position-error feedback row
($[\insF]_{\delta\rxpos,\delta\rxpos} = -\skewm{\earthrate}$) and the
gravity-gradient coupling
$[\insF]_{\delta\imuvel,\delta\rxpos} = +\,\partial\gravec/\partial\rxpos$
(the gradient of the \emph{plumb} gravity of \cref{eq:imu-gravity}, which
already subsumes the centrifugal $-\skewm{\earthrate}^{2}$; the leading
sign follows from differentiating the velocity update
\cref{eq:imu-velocity}).  The resulting first-order model is accurate
over the short epoch interval of a GNSS-aided run.  The discrete
transition is also first order,
\begin{equation}
[\insPhi]_k = \boldsymbol I + \insF\,\dt .
\label{eq:ins-Phi}
\end{equation}
The process noise maps the continuous accelerometer and gyro white
noise and the bias random walks through
\begin{equation}
[\insQ]_k = \insG\,\boldsymbol W\,\insG^{\Tr}\,\dt, \quad
\insG =
\begin{bmatrix}
\boldsymbol 0 & \boldsymbol 0 & \boldsymbol 0 & \boldsymbol 0\\
\Cbe & \boldsymbol 0 & \boldsymbol 0 & \boldsymbol 0\\
\boldsymbol 0 & -\Cbe & \boldsymbol 0 & \boldsymbol 0\\
\boldsymbol 0 & \boldsymbol 0 & \boldsymbol I & \boldsymbol 0\\
\boldsymbol 0 & \boldsymbol 0 & \boldsymbol 0 & \boldsymbol I
\end{bmatrix}, \quad
\boldsymbol W = \operatorname{diag}\!\bigl(
  \sigma_a^2, \sigma_g^2, \sigma_{ba}^2, \sigma_{bg}^2\bigr)\!\otimes\! \boldsymbol I_3 ,
\label{eq:ins-Q}
\end{equation}
where the spectral densities $\sigma_a,\sigma_g,\sigma_{ba},\sigma_{bg}$
are read from
\texttt{imu.\{accel\_noise\_density},
\texttt{gyro\_noise\_density}, \texttt{accel\_bias\_instability},
\texttt{gyro\_bias\_instability\}}.  The defaults in
\texttt{config/GnssConfigTypes.hpp} are
\SI{0.05}{\metre\per\second\squared\per\hertz\tothe{0.5}},
\SI{5e-4}{\radian\per\second\per\hertz\tothe{0.5}},
\SI{1e-4}{\metre\per\second\cubed\per\hertz\tothe{0.5}} and
\SI{1e-6}{\radian\per\second\squared\per\hertz\tothe{0.5}}, in that
order.  Over one GNSS epoch the
per-sample pairs $([\insPhi]_k,[\insQ]_k)$ accumulate by the standard
recursion
\begin{equation}
\insPhi \leftarrow [\insPhi]_k\,\insPhi, \qquad
\insQ   \leftarrow [\insPhi]_k\,\insQ\,[\insPhi]_k^{\Tr} + [\insQ]_k .
\label{eq:ins-accum}
\end{equation}

\subsection{Covariance splice and attitude feedback}\label{sec:imu-splice}

The accumulated $15\times15$ pair $(\insPhi,\insQ)$ is written over the
indices the INS owns in an identity-seeded full-state transition $\Fmat$
and in the random-walk process-noise baseline $\Qmat$.  The write is an
assignment, so the baseline noise on those rows is not counted twice.
The baseline process noise is evaluated over the interval the INS
integrated.  The non-INS states (clock, troposphere, ionosphere, ambiguities, \dots)
keep their random-walk model from \dcref{sec:kf-predict}{the time
update}; because the transition is the identity there, their means are
not propagated.  The covariance is then propagated once and symmetrised,
\begin{empheq}[box=\fbox]{equation}
\covM \leftarrow \Fmat\,\covM\,\Fmat^{\Tr} + \Qmat, \qquad
\covM \leftarrow \tfrac12\bigl(\covM + \covM^{\Tr}\bigr).
\label{eq:imu-cov-prop}
\end{empheq}
The mechanised values of \crefrange{eq:imu-velocity}{eq:imu-position}
overwrite the nominal navigation states $(\rxpos,\imuvel)$.  The
attitude error is reset, $\imupsi := \boldsymbol 0$, because
\cref{eq:imu-attitude} has already applied the rotation increment to
$\Cbe$: after each aided prediction the IMU aiding source writes the
propagated attitude into the filter through its attitude setter.  The
filter holds the single nominal attitude.  It is initialised to the
identity (there is no initial alignment step), and on the live path only
two operations write it: that setter call and the feedback step below.
After the next GNSS measurement update, the corrected attitude error
feeds back multiplicatively into the DCM and is reset again,
\begin{equation}
\Cbe \leftarrow \bigl(\boldsymbol I + \skewm{\imupsi}\bigr)\Cbe,
\qquad \imupsi := \boldsymbol 0 ,
\label{eq:imu-att-feedback}
\end{equation}
followed by SVD re-orthonormalisation
(\texttt{KalmanFilter::update}).  This closed-loop reset keeps $\imupsi$
a small angle, so the linearisation in \cref{eq:ins-F} remains valid.

\subsection{Coupling modes and the GNSS update}\label{sec:imu-coupling}

The coupling mode (\texttt{imu.coupling\_mode}) determines how the GNSS
observations correct the inertial errors.  In both modes the filter
processes the raw GNSS observations; the modes differ in the Doppler
Jacobian and in the relinearisation path (see the remark at the end of
this section).

\paragraph{Loosely coupled (default).}
Among the INS states, the GNSS Jacobians of
\dcref{sec:obs-geometry}{the observation geometry} touch only the
position and velocity columns.  The INS error states
$\imupsi,\biasa,\biasg$ have no direct observation partials.  They are
corrected indirectly, through the cross-covariances that
\cref{eq:imu-cov-prop} builds between the position/velocity block and the
attitude/bias block.  The feedback of \cref{eq:imu-att-feedback} then
applies the attitude correction.

\paragraph{Tightly coupled.}
Tight coupling requires the velocity block, so the loader enables it
regardless of \texttt{estimation.estimate\_velocity}.  A
\texttt{TightCoupled} \texttt{coupling\_mode} with the flag false sets
the flag true and logs \texttt{"forcing estimate\_velocity=true"}
(\texttt{GnssConfigLoad.cpp}).  The three velocity states, and with
them the Doppler rows that carry the Jacobian below, are therefore
always present.  In this mode the Doppler observation, and only the
Doppler observation, gains direct attitude and accelerometer-bias
columns.  It is the one aiding contribution that adds columns to
existing GNSS rows instead of appending rows of its own.  Each
relinearisation pass calls
\texttt{AidingSourceManager::applyTcJacobianInputs} $\to$
\texttt{ImuAidingSource::applyTcJacobianInputs}, which passes the current
attitude and specific force to \texttt{MeasurementModel::setImuState}
(\texttt{src/core/measurement/MeasurementModel.cpp}).  With
$\imufe = \Cbe\imuf$ the specific force in ECEF and
$\los^\sat_{\rx}$ the unit line-of-sight of
\dcref{eq:los}{the line-of-sight definition}, the Doppler measurement
function $h_{\mathrm{D}}$
contributes
\begin{empheq}[box=\fbox]{equation}
\frac{\partial h_{\mathrm{D}}}{\partial \imupsi}
  = -\,\skewm{\imufe}^{\Tr}\,\los^\sat_{\rx}, \qquad
\frac{\partial h_{\mathrm{D}}}{\partial \biasa}
  = -\,(\Cbe)^{\Tr}\,\los^\sat_{\rx}.
\label{eq:imu-tc-jac}
\end{empheq}
The specific force used in \cref{eq:imu-tc-jac} is the static-rest
approximation $\imuf \approx (\Cbe)^{\Tr}\!\bigl(-\gravec\bigr) - \biasa$
(\texttt{ImuAidingSource.cpp}), evaluated at the current position
estimate rather than from the IMU samples.  It is exact only when the
true platform acceleration is negligible compared with gravity.  In
either mode the code and carrier-phase rows carry no attitude or bias
partials; only the Doppler rows couple directly to the inertial errors.

\begin{remark}[Relinearisation under tight coupling]
The tightly coupled path keeps a single \emph{naive} second-pass
relinearisation instead of the prior-anchored iterated EKF of
\dcref{sec:kf-iekf}{the iterated EKF}.  Re-deriving the gain from a
fixed prior is inconsistent with the multiplicative attitude feedback
of \cref{eq:imu-att-feedback}, which has already changed $\Cbe$ and
zeroed the attitude error between passes.  The aiding layer exposes the
requirement as a capability query
(\texttt{AidingSourceManager::anySourceRequiresNaiveRelin\-earization},
true for the IMU source under tight coupling), so any later source with
nominal-state feedback selects the same path.  Direct
\texttt{tightlyCoupledImu()} checks remain in
\texttt{Application.cpp} / \texttt{KalmanFilter.cpp} (the IEKF gating
and the prior capture), so the migration to the capability query is
incomplete.
\end{remark}

\ifSubfilesClassLoaded{\printbibliography}{}
\end{document}
